\appsection{Ergänzende Trainingsverläufe Seed 0}
\label{app:training_curves}

Alle Verläufe sind mit exponentiellem gleitendem Mittelwert (EMA, $\alpha = 0{,}9$) geglättet.

\begin{figure}[H]
    \centering
    \includegraphics[width=0.48\textwidth]{figures/training/Episode_rew_progress_raw.pdf}
    \hfill
    \includegraphics[width=0.48\textwidth]{figures/training/Episode_rew_close_raw.pdf}
    \caption[Fortschrittsbelohnung und Näherungsbelohnung]{Fortschrittsbelohnung
    (links) und Näherungsbelohnung (rechts). Beide Komponenten steigen mit
    zunehmender Annäherung an das Ziel.}
    \label{fig:rew_progress_close}
\end{figure}

\begin{figure}[H]
    \centering
    \includegraphics[width=0.48\textwidth]{figures/training/Episode_rew_obstacle_proximity_raw.pdf}
    \hfill
    \includegraphics[width=0.48\textwidth]{figures/training/Episode_rew_crash_raw.pdf}
    \caption[Hindernisabstandsstrafe und Absturzstrafe]{Hindernisabstandsstrafe
    (links) und Absturzstrafe (rechts). Beide Strafkomponenten fallen mit
    zunehmender Kollisionsvermeidung.}
    \label{fig:rew_prox_crash}
\end{figure}

\begin{figure}[H]
    \centering
    \includegraphics[width=0.48\textwidth]{figures/training/Episode_rew_success_raw.pdf}
    \hfill
    \includegraphics[width=0.48\textwidth]{figures/training/Train_mean_episode_length.pdf}
    \caption[Erfolgsbelohnung und Episodenlänge]{Erfolgsbelohnung (links) und
    mittlere Episodenlänge (rechts). Der Anstieg der Erfolgsbelohnung markiert
    den Beginn erfolgreicher Landungen. Die Episodenlänge gibt die Anzahl der
    Entscheidungsschritte pro Episode an.}
    \label{fig:rew_success_eplen}
\end{figure}

\begin{figure}[H]
    \centering
    \includegraphics[width=0.48\textwidth]{figures/training/Policy_mean_std.pdf}
    \hfill
    \includegraphics[width=0.48\textwidth]{figures/training/Loss_entropy.pdf}
    \caption[Standardabweichung und Entropie]{Mittlere Standardabweichung der
    Aktionsverteilung (links) und Entropie (rechts). Beide Maße fallen
    monoton und spiegeln die abnehmende Exploration wider.}
    \label{fig:rew_prox_std}
\end{figure}

\begin{figure}[H]
    \centering
    \includegraphics[width=0.48\textwidth]{figures/training/Loss_surrogate.pdf}
    \hfill
    \includegraphics[width=0.48\textwidth]{figures/training/Loss_value.pdf}
    \caption[Surrogatverlust und Wertverlust]{PPO-Surrogatverlust (links) und
    Wertverlust (rechts).}
    \label{fig:loss_surrogate_value}
\end{figure}


\appsection{Standardabweichung über alle Seeds}
\label{app:std_seeds}

\begin{figure}[H]
    \centering
    \includegraphics[width=\textwidth]{Bilder/seed_analysis/policy_std_seeds.pdf}
    \caption[Standardabweichung der Aktionsverteilung]{Mittlere
    Standardabweichung der Aktionsverteilung über alle sechs Seeds. Der
    Verlauf ist nahezu identisch und fällt von 0,99 auf circa 0,15,
    unabhängig vom jeweiligen Belohnungsverlauf. Geglättet mit EMA
    $\alpha = 0{,}9$.}
    \label{fig:policy_std_seeds}
\end{figure}


\appsection{Phase-1-Evaluationsmetriken}
\label{app:phase1_metrics}

\begin{table}[H]
  \centering
  \caption[Phase-1-Ergebnisse des RL-Agenten]{Phase-1-Ergebnisse des
  RL-Agenten ($1\,000$~Episoden, Checkpoint~400, ungesehene Hindernisse).
  Konfidenzintervalle: 95\,\%-Bootstrap (Perzentilmethode, $B = 10\,000$).
  Landezeit und Pfadeffizienz beziehen sich ausschließlich auf
  erfolgreiche Episoden ($n = 710$).}
  \label{tab:phase1_results}
  \small
  \renewcommand{\arraystretch}{1.1}
  \begin{tabular}{l r r}
    \toprule
    \textbf{Metrik} & \textbf{Wert} & \textbf{95\,\%-KI / IQR} \\
    \midrule
    Erfolgsrate                       & 71,0\,\%          & [68,2;\; 73,8] \\
    \addlinespace
    Hinderniskollision                & 28,8\,\% (288)    & [26,0;\; 31,6] \\
    Verlassen des Korridors           & 0,2\,\% (2)       & [0,0;\; 0,5] \\
    Bodenkollision                    & 0,0\,\% (0)       & — \\
    Timeout                           & 0,0\,\% (0)       & — \\
    \addlinespace
    Landezeit (Median)                & 30,0\,s            & IQR [27,0;\; 36,0] \\
    Pfadeffizienz (Median)            & 1,19               & IQR [1,13;\; 1,25] \\
    Min.\ Hindernisabstand (Mittel)   & 0,52\,m            & [0,50;\; 0,54] \\
    \bottomrule
  \end{tabular}
\end{table}


\appsection{Phase-1-Evaluationsmetriken (oberflächenbasierte Kollisionserkennung)}
\label{app:phase1_mesh_metrics}

\begin{table}[H]
  \centering
  \caption[Phase-1-Ergebnisse mit oberflächenbasierter Kollisionserkennung]{Phase-1-Ergebnisse des
  RL-Agenten mit oberflächenbasierter Kollisionserkennung ($1\,000$~Episoden,
  Checkpoint~400, ungesehene Hindernisse).
  Konfidenzintervalle: 95\,\%-Bootstrap (Perzentilmethode, $B = 10\,000$).
  Landezeit und Pfadeffizienz beziehen sich ausschließlich auf
  erfolgreiche Episoden ($n = 984$).}
  \label{tab:phase1_mesh_results}
  \small
  \renewcommand{\arraystretch}{1.1}
  \begin{tabular}{l r r}
    \toprule
    \textbf{Metrik} & \textbf{Wert} & \textbf{95\,\%-KI / IQR} \\
    \midrule
    Erfolgsrate                       & 98,4\,\%          & [97,6;\; 99,1] \\
    \addlinespace
    Hinderniskollision                & 1,2\,\% (12)      & [0,6;\; 1,9] \\
    Verlassen des Korridors           & 0,2\,\% (2)       & [0,0;\; 0,5] \\
    Bodenkollision                    & 0,0\,\% (0)       & — \\
    Timeout                           & 0,2\,\% (2)       & [0,0;\; 0,5] \\
    \addlinespace
    Landezeit (Median)                & 33,0\,s            & IQR [30,0;\; 36,0] \\
    Pfadeffizienz (Median)            & 1,19               & IQR [1,13;\; 1,25] \\
    Min.\ Hindernisabstand (Mittel)   & 0,51\,m            & [0,49;\; 0,53] \\
    \bottomrule
  \end{tabular}
\end{table}


\appsection{Phase-1-Evaluationsmetriken (RRT-Connect)}
\label{app:phase1_rrt_metrics}

\begin{table}[H]
  \centering
  \caption[Phase-1-Ergebnisse des RRT-Connect-Planers]{Phase-1-Ergebnisse des
  RRT-Connect-Planers ($1\,000$~Episoden, Checkpoint~400, ungesehene Hindernisse).
  Konfidenzintervalle: 95\,\%-Bootstrap (Perzentilmethode, $B = 10\,000$).
  Landezeit und Pfadeffizienz beziehen sich ausschließlich auf
  erfolgreiche Episoden ($n = 820$).}
  \label{tab:phase1_rrt_results}
  \small
  \renewcommand{\arraystretch}{1.1}
  \begin{tabular}{l r r}
    \toprule
    \textbf{Metrik} & \textbf{Wert} & \textbf{95\,\%-KI / IQR} \\
    \midrule
    Erfolgsrate                       & 82,0\,\%          & [79,6;\; 84,4] \\
    \addlinespace
    Hinderniskollision                & 18,0\,\% (180)    & [15,6;\; 20,4] \\
    Verlassen des Korridors           & 0,0\,\% (0)       & — \\
    Bodenkollision                    & 0,0\,\% (0)       & — \\
    Timeout                           & 0,0\,\% (0)       & — \\
    \addlinespace
    Landezeit (Median)                & 42,0\,s            & IQR [42,0;\; 48,0] \\
    Pfadeffizienz (Median)            & 1,09               & IQR [1,07;\; 1,20] \\
    Min.\ Hindernisabstand (Mittel)   & 0,48\,m            & [0,47;\; 0,49] \\
    \bottomrule
  \end{tabular}
\end{table}


\appsection{Phase-2-Evaluationsmetriken}
\label{app:phase2_metrics}

\begin{table}[H]
  \centering
  \caption[Phase-2-Ergebnisse des RL-Agenten]{Phase-2-Ergebnisse des
  RL-Agenten ($1\,000$~Episoden, Checkpoint~400, Weinberg-Transfer,
  Erfolgsradius $0{,}5$\,m).
  Konfidenzintervalle: 95\,\%-Bootstrap (Perzentilmethode, $B = 10\,000$).
  Landezeit und Pfadeffizienz beziehen sich ausschließlich auf
  erfolgreiche Episoden ($n = 104$).}
  \label{tab:phase2_results}
  \small
  \renewcommand{\arraystretch}{1.1}
  \begin{tabular}{l r r}
    \toprule
    \textbf{Metrik} & \textbf{Wert} & \textbf{95\,\%-KI / IQR} \\
    \midrule
    Erfolgsrate                       & 10,4\,\%          & [8,5;\; 12,4] \\
    \addlinespace
    Terrain-Kollision                 & 15,9\,\% (159)    & [13,6;\; 18,2] \\
    Timeout                           & 73,7\,\% (737)    & [71,0;\; 76,4] \\
    Bodenkollision                    & \multicolumn{2}{c}{entfällt (kein fester Bodenkollisions-Check)} \\
    Verlassen des Korridors           & \multicolumn{2}{c}{entfällt (kein Korridorbegrenzungsrahmen)} \\
    \addlinespace
    Landezeit (Median)                & 36,0\,s            & IQR [33,0;\; 36,0] \\
    Pfadeffizienz (Median)            & 1,07               & IQR [1,05;\; 1,09] \\
    Finale Zieldistanz (Median)       & 1,19\,m            & — \\
    Min.\ Terrain-Abstand (Mittel)    & 1,02\,m            & [0,98;\; 1,06] \\
    \bottomrule
  \end{tabular}
\end{table}


\appsection{Phase-2-Evaluationsmetriken (RRT-Connect)}
\label{app:phase2_rrt_metrics}

\begin{table}[H]
  \centering
  \caption[Phase-2-Ergebnisse des RRT-Connect-Planers]{Phase-2-Ergebnisse des
  RRT-Connect-Planers ($1\,000$~Episoden, Checkpoint~400, Weinberg-Transfer,
  Erfolgsradius $0{,}5$\,m).
  Konfidenzintervalle: 95\,\%-Bootstrap (Perzentilmethode, $B = 10\,000$).
  Landezeit und Pfadeffizienz beziehen sich ausschließlich auf
  erfolgreiche Episoden ($n = 857$).}
  \label{tab:phase2_rrt_results}
  \small
  \renewcommand{\arraystretch}{1.1}
  \begin{tabular}{l r r}
    \toprule
    \textbf{Metrik} & \textbf{Wert} & \textbf{95\,\%-KI / IQR} \\
    \midrule
    Erfolgsrate                       & 85,7\,\%          & [83,5;\; 87,9] \\
    \addlinespace
    Terrain-Kollision                 & 7,1\,\% (71)      & [5,6;\; 8,7] \\
    Timeout                           & 7,9\,\% (79)      & [6,3;\; 9,6] \\
    Bodenkollision                    & \multicolumn{2}{c}{entfällt (kein fester Bodenkollisions-Check)} \\
    Verlassen des Korridors           & \multicolumn{2}{c}{entfällt (kein Korridorbegrenzungsrahmen)} \\
    \addlinespace
    Landezeit (Median)                & 42,0\,s            & IQR [39,0;\; 48,0] \\
    Pfadeffizienz (Median)            & 0,99               & IQR [0,97;\; 1,00] \\
    Min.\ Terrain-Abstand (Mittel)    & 0,46\,m            & [0,45;\; 0,47] \\
    \bottomrule
  \end{tabular}
\end{table}


\appsection{PID-Reglerparameter}
\label{app:pid}

\begin{table}[H]
  \centering
  \caption{PID-Verstärkungen des kaskadierenden Reglers.}
  \label{tab:app_pid_gains}
  \begin{tabular}{l c c c}
    \toprule
    Regelkreis & $K_p$ & $K_i$ & $K_d$ \\
    \midrule
    Position $x$/$y$         & 1{,}0   & 0     & 0{,}7  \\
    Position $z$             & 1{,}5   & 0     & 1{,}0  \\
    Geschwindigkeit $x$/$y$  & 16{,}0  & 0     & 8{,}0  \\
    Geschwindigkeit $z$      & 100{,}0 & 2{,}0 & 10{,}0 \\
    Roll / Nick              & 6{,}0   & 0     & 3{,}0  \\
    Gier                     & 0{,}5   & 0     & 0{,}8  \\
    \bottomrule
  \end{tabular}
\end{table}
